\documentclass[11pt,a4paper]{article}
\usepackage[utf8]{inputenc}
\usepackage[top=0.9in, bottom=0.8in, left=1in, right=1in]{geometry}
\usepackage{hyperref}

\setlength{\parskip}{0.8em}
\setlength{\parindent}{0pt}

\begin{document}

\begin{center}
    {\LARGE \textbf{Liuxi Mei}} \\
    \vspace{2mm}
    Linköping, Sweden $\cdot$ liume102@student.liu.se $\cdot$ +46-0704898962
    \vspace{4mm}
\end{center}

\begin{flushright}
April 9, 2026
\end{flushright}

Application for the doctoral student position in Machine Learning


I am writing to apply for the PhD student position in Machine Learning at Linköping University. I am finishing my master's degree here now. I will graduate in June this year.

\textbf{Motivation for the position:} I really like this position because it goes beyond normal wireless communication. Normal communication is mostly about sending data. But this project is about sensing and understanding the real world. Also, combining telecommunications with artificial intelligence is very exciting to me. In the age of AI, having exclusive data is extremely important for building good models. This project features a great lab setup. This means I can use unique, real-world data instead of just computer simulations. Finally, the academic environment and the lovely city make Linköping University the perfect place for my PhD.

\textbf{Professional and academic background:} My background brings together telecommunications, software engineering, and machine learning. Before my master's degree, I worked for over five years as a software engineer at Ericsson. I built an energy-saving platform for telecom networks. This platform turned off MIMO components to save power. I know this is different from the D-MIMO focus of your project. But this work gave me good experience with telecom hardware. It also taught me to care about system reliability and fast code. During my master's degree, I moved closer to research. I learned to work with large datasets and build machine learning models in Python.

\textbf{Master's thesis research:} My master's thesis focuses on graph neural networks. I developed deep learning models to analyze complex, structured data. While the application area is not wireless communication, the core research process is highly relevant. I have gained strong experience in evaluating model performance and understanding why certain architectures succeed. I am also highly skilled at designing rigorous experiments and handling noisy, real-world datasets.

\textbf{Suitability for research topic:} My background combines machine learning, software engineering, and telecommunications. My master's thesis on graph neural networks (GNNs) provided direct experience in building models for complex network structures. This GNN expertise is highly applicable to this project because a distributed MIMO system can be naturally modeled as a physical graph. Furthermore, my five years of engineering work at Ericsson built a practical understanding of telecom hardware and network traffic. I believe this mix of graph-based machine learning and industry telecommunication experience is a strong fit for the demands of this research.

Thank you for considering my application. I would be very happy to contribute to the research environment at STIMA and to this project.

Sincerely,

Liuxi Mei

\vspace{0.5cm}
\noindent
\begin{minipage}{\textwidth}
\textbf{References:}

\vspace{0.3cm}
\begin{minipage}[t]{0.48\textwidth}
\textbf{Oleg Sysoev}\\
Senior Associate Professor, Linköping University\\
Department of Computer and Information Science (IDA)\\
Master's Thesis Supervisor\\
Email: \href{mailto:oleg.sysoev@liu.se}{oleg.sysoev@liu.se}
\end{minipage}%
\hfill
\begin{minipage}[t]{0.48\textwidth}
\textbf{Krzysztof Bartoszek}\\
Senior Associate Professor, Linköping University\\
Department of Computer and Information Science (IDA)\\
Teacher and Examiner\\
Email: \href{mailto:krzysztof.bartoszek@liu.se}{krzysztof.bartoszek@liu.se}
\end{minipage}
\end{minipage}

\end{document}
